\documentclass[11pt,a4paper]{article}
\usepackage[T1]{fontenc}
\usepackage{isabelle,isabellesym}

\usepackage{amssymb}

% this should be the last package used
\usepackage{pdfsetup}

% urls in roman style, theory text in italics
\urlstyle{rm}
\isabellestyle{it}

\begin{document}

\title{Ownership: an Instantiation of the Generic Owicki--Gries Logic, with
  Use-After-Free Safety}
\author{Christian Kissig}
\maketitle

\begin{abstract}
  This entry is a concrete instantiation of the generic Owicki--Gries logic
  (AFP entry \emph{OwickiGries}).  Actions are tagged with the acting process,
  resources have an abstract type, and the abstract state carries an
  \emph{ownership map} assigning a set of owned resources to each process.  A
  three-rule transfer discipline (locality, footprint, and fresh acquisition)
  fixes how actions move ownership; it abstracts the rules of a read-copy-update
  (RCU) algorithm --- a thread assumes ownership of freshly allocated resources,
  gains ownership after reading the shared structure, and relinquishes it at the
  rcu-exit / \texttt{free}.

  From the discipline we prove an \emph{exclusivity} invariant (no resource is
  owned by two processes) and its preservation under reachability.  The central
  reusable result is that \emph{ownership disjointness discharges
  interference-freedom}: because a foreign action leaves a process's ownership
  untouched, every assertion tagged to another thread is automatically stable,
  so the Owicki--Gries interference-freedom obligation holds by construction.
  Finally we derive a \emph{use-after-free safety} corollary: the configuration
  that co-enables a \texttt{free} with a still-pending dereference of the same
  resource is unreachable, so the logic is stuck before witnessing the bug.

  The entry depends only on its sibling \emph{OwickiGries} (no further imports).
\end{abstract}

\tableofcontents

% sane default for proof documents
\parindent 0pt\parskip 0.5ex

% generated text of all theories
\input{session}

\end{document}
